\section{Definitions and Foundations}

Before delving into a comprehensive survey, we first present a formal definition of self-evolving agents and introduce a taxonomy of the key aspects in self-evolving agents. We also discuss the relationships between self-evolving agents and other renowned learning paradigms, such as curriculum learning, lifelong learning, model editing, and unlearning, highlighting the adaptive, dynamic, and autonomous nature of self-evolving agents.

\subsection{Definitions}
\label{sec:definition}

\paragraph{Environment} We first define the environment (including the user and the execution environment, e.g., Linux shell) of an agent system as a partially observable Markov Decision Process (POMDP), represented as a tuple $E=(\mathcal{G}, \mathcal{S}, \mathcal{A}, T, R, \Omega, O, \gamma)$, where:
\begin{itemize}
\item $\mathcal{G}$ is a set of potential goals. Each $g\in \mathcal{G}$ is a task objective that the agent needs to achieve, e.g., a user query.
\item $\mathcal{S}$ is a set of states. Each $s\in \mathcal{S}$ represents the internal state of the environment.
\item $\mathcal{A}$ is a set of actions. Each action $a\in \mathcal{A}$ can be a combination of textual reasoning, retrieval of external knowledge, and tool calls.
\item $T$ is the state transition probability function which takes a state-action pair $(s,a)$ and outputs the probability distribution $T(s'|s,a)$ of the next state.
\item $R:\mathcal{S}\times \mathcal{A} \times \mathcal{G} \rightarrow \mathcal{R}$ is the feedback/reward function, conditioned on the specific goal $g\in \mathcal{G}$. The feedback $r=R(s,a,g)$ typically takes the form of a scalar score or textual feedback.
\item $\Omega$ is a set of observations accessible to the agent.
\item $O$ is the observation probability function which takes a state-action pair $(s,a)$ and outputs the probability distribution $O(o'|s,a)$ of the next observation for the agent.
\item $\gamma$ is the discount factor.
\end{itemize}

\paragraph{Agent system} We define a (multi-)agent system as $\Pi =(\Gamma,\{\psi_i\}, \{C_i\},\{\mathcal{W}_i\})$. The architecture $\Gamma$ determines the control flow of the agent system or collaborative structures between multiple agents. It is typically represented as a sequence of nodes $(N_1, N_2, ...)$ organized by graph or code structures. Each node $N_i$ consists of the following components:
\begin{itemize}
\item $\psi_i$: the underlying LLM/MLLM.
\item $C_i$: the context information, e.g., prompt $P_i$ and memory $M_i$.
\item $\mathcal{W}_i$: the set of available tools/APIs.
\end{itemize}

At each node, the agent policy is a function $\pi_{\theta_i}(\cdot|o)$ that takes an observation and outputs the probability distribution of the next action, where $\theta_i=(\psi_i, C_i)$. The actual action space here is the union of the natural language space and the tool space $\mathcal{W}_i$.

For a given task $\mathcal{T}=(E, g)$, represented by an environment $E$ and a corresponding goal $g\in \mathcal{G}$, the agent system follows the topology $\Gamma$ to generate a trajectory $\tau=(o_0, a_0, o_1, a_1,...)$, and receives a feedback $r$ either from the external environment or from internal signals (e.g., self-confidence or feedback from an evaluator).

\paragraph{Self-evolving strategy} A self-evolving strategy is a transformation $f$ that maps the current agent system to a new state, conditioned on the generated trajectory $\tau$ and the external/internal feedback $r$:  
\begin{equation}
    f(\Pi, \tau, r)=\Pi'=(\Gamma',\{\psi'_i\}, \{C'_i\}, \{\mathcal{W}'_i\})
\end{equation}

\paragraph{Objective of self-evolving agents} 
Let $U$ be a utility function that measures the performance of an agent system $\Pi$ on a given task $\mathcal{T}$ by assigning a scalar score $U(\Pi, \mathcal{T}) \in \mathbb{R}$. The utility may be derived from the task-specific feedback $r$, such as a reward signal or textual evaluation, possibly combined with other performance indicators (e.g., completion time, accuracy, or robustness). Given a sequence of tasks $(\mathcal{T}_0, \mathcal{T}_1, ..., \mathcal{T}_n)$ and an initial agent system $\Pi_0$, a self-evolving strategy $f$ recurrently generates an evolving sequence of agent systems $(\Pi_1, \Pi_2, ..., \Pi_n)$ via
\begin{equation}
    \Pi_{j+1} = f(\Pi_j, \tau_j, r_j),
\end{equation}
where $\tau_j$ and $r_j$ are the trajectory and feedback on task $\mathcal{T}_j$. 

The overarching objective in designing a self-evolving agent is to construct a strategy $f$ such that the cumulative utility over tasks is maximized:
\begin{equation}
    \max_f \sum_{j=0}^n U(\Pi_j, \mathcal{T}_j)
\end{equation}

\paragraph{Operational definition of self-evolving agents}
To provide a conceptual boundary, we introduce an operational definition of self-evolving agents. 
A self-evolving agent is the agent that \emph{modifies its internal parameters, contextual state, toolset, or architectural topology based on its own trajectories or feedback signals, with the explicit objective of improving future performance}. 

This definition entails three inclusion criteria: 
(i) updates must be \emph{experience-dependent}, driven by trajectories, self-generated data, or environment feedback, specifically targeting the agent's policy limitations or capability boundaries rather than generic data synthesis; 
(ii) updates must produce a \emph{persistent, policy-changing} effect rather than a transient instruction-following behavior; 
(iii) the system must possess mechanisms for \emph{autonomous exploration or self-initiated learning}, even if it also leverages pre-collected data. For clarity, we use "passive" to denote learning triggered exclusively by externally provided data or schedules, and "active" to denote self-initiated exploration, reflection, or structural modification (i.e., using self-reflection to collect data), explicitly excluding static pipelines (e.g., standard distillation) where data generation is agnostic to the agent's interaction history.

As this field is rapidly forming, fully autonomous self-evolution without human intervention represents an aspirational goal rather than the current norm. In this survey, we do not impose a rigid exclusion threshold that would disregard early-stage developments. Instead, we analyze the mechanisms contributing to the self-evolving paradigm ranging from \textbf{proto-evolution} (e.g., iterative bootstrapping or feedback-driven prompting) to \textbf{strong self-evolution} (fully autonomous diagnosis and reconfiguration), allowing us to provide a comprehensive view of how diverse methods contribute to the "What, When, and How" of the paradigm's progression toward full autonomy.





\subsection{Relationships with Other Works}


Table \ref{tab:compare-paradigm} summarizes the key distinctions between self-evolving agents and other paradigms (including curriculum learning, lifelong learning, model editing, and unlearning). 
We provide a brief introduction to each paradigm below,  highlighting the differences among these paradigms, as well as the differences with self-evolving agents.

\paragraph{Curriculum Learning}
Curriculum learning is a training strategy in which data are presented in order of increasing difficulty \citep{bengio2009curriculum, wang2021survey}. This strategy resembles human curricula where concepts are introduced progressively from simple to complex. Curriculum learning has been widely adopted across diverse domains, including computer vision \citep{guo2018curriculumnet, jiang2014easy, liu2023towards}, natural language processing \citep{platanios2019competence, tay2019simple}, speech recognition \citep{braun2017curriculum, lotfian2019curriculum}, etc. Recently, several curriculum learning-based methods have been proposed to fine-tune LLMs during the post-training phase \citep{wang2025dump, zhang2025preference, parashar2025curriculum, zhang2025learning, li2025curriculum}. The framework for curriculum learning generally comprises two key components: a difficulty measurer that quantifies the difficulty level of each training data point, and a training scheduler that reorganizes the order of data points received by the model according to the difficulty level. 
Unlike curriculum learning, which operates on a static dataset, self-evolving agents aim to handle sequential tasks in dynamic environments. Additionally, curriculum learning updates only model parameters, whereas self-evolving agents are able to adjust non-parametric components like memory and tools.

\paragraph{Lifelong Learning}
Lifelong learning refers to the ability of AI models to continuously and adaptively learn when exposed to new tasks and environments, while retaining previously acquired knowledge and abilities. This learning paradigm, also known as continual learning or incremental learning, is crucial for AI models to operate in dynamic and complex environments \citep{wang2024comprehensive, zheng2025lifelong, parisi2019continual, shi2024continual, yang2025recent, zhou2024continual}. The primary goal of lifelong learning for AI models is to achieve a balance between preserving existing knowledge (stability) and acquiring new knowledge (plasticity) when exposed to new data or tasks \citep{mccloskey1989catastrophic, zheng2025lifelong, ratcliff1990connectionist, rolnick2019experience}. Though it shares the sequential task setting with self-evolving agents, lifelong learning differs in two fundamental ways: (1) \textit{Memory functionality and usage timing}: While continual learning methods extensively employ memory mechanisms (e.g., experience replay buffers \citep{rolnick2019experience}, episodic memory \citep{lopez2017gradient}) to mitigate catastrophic forgetting, these mechanisms primarily serve as \emph{training-time} tools for parameter optimization through gradient computation. In contrast, self-evolving agents leverage \emph{runtime context} (prompts, working memory, conversation history) that directly influences action generation at test-time without requiring parameter updates. The distinction lies not in the presence of non-parametric components, but in their functional role: training-time replay vs. test-time state adaptation. (2) \textit{Learning initiative}: Lifelong learning primarily acquires knowledge passively through externally provided task sequences, whereas self-evolving agents actively explore their environment and incorporate internal reflection or self-evaluation mechanisms to guide their own learning trajectory. Recent self-improving LLM methods \citep{huang2022large, yuan2024self}, which iteratively refine models through self-generated data and self-critique, can be viewed as instances of lifelong learning focused on model-centric improvement. Self-evolving agents extend beyond this paradigm to encompass system-wide evolution including tool acquisition, architectural reconfiguration, and environmental exploration.

\paragraph{Model Editing and Unlearning}
Model editing and unlearning aim to efficiently and precisely modify specific knowledge in AI models while preserving irrelevant knowledge and avoiding full retraining \citep{wang2024knowledge, wang2025comprehensive, zhang2024comprehensive, wang2025comprehensive, nguyen2022survey, geng2025comprehensive}. A canonical application of model editing is to perform efficient and precise localized factual updates (e.g., modifying the answer to "2021 Olympics host city" from "Tokyo" to "Paris"). Early methods focused on triples of atomic knowledge and later expanded into various trustworthy-related tasks \citep{fang2025alphaedit, huang2025model}. Recent studies also propose lifelong model editing\citep{chen2024lifelong} that sequentially performs model editing.
For model unlearning, early efforts mainly focus on the removal of privacy-related information \citep{chen2021machine}. With the rapid development of LLMs, model unlearning is also used to enhance LLMs' safety \citep{safe_unlearning, li2024wmdp, circuit_breaker, lu2025x}. 
Compared to lifelong learning, model editing shares an aligned objective: both aim to acquire new knowledge or capabilities while mitigating catastrophic forgetting. However, lifelong learning typically relies on extensive gradient-based fine-tuning across all model parameters, whereas model editing often modifies only a small subset of parameters in a targeted manner.
Compared to self-evolving agents, model editing (1) cannot modify non-parametric components such as memory or tools, and (2) relies on a pre-defined pipeline from the algorithm designer, whereas self-evolving agents can spontaneously employ more diverse and flexible strategies based on the observation of the environment or internal feedback signals.

\begin{table}[t]
\small
\centering
\caption{Comparison between self-evolving agents and other renowned paradigms}
\resizebox{\textwidth}{!}{
\begin{tabular}{@{}lccccccc@{}}
\toprule
\multirow{2}{*}{\textbf{Paradigm}} & \textbf{Runtime} & \textbf{Evolving} & \textbf{Dynamic} & \textbf{Test-time} & \textbf{Active} & \textbf{Structural} & \textbf{Self-reflect} \\
& \textbf{Context} & \textbf{Toolset} & \textbf{Tasks} &\textbf{Adaptation} & \textbf{Exploration} & \textbf{Change} & \textbf{\& Eval} \\
\midrule
Curriculum Learning & \textcolor{myred}{\ding{55}} & \textcolor{myred}{\ding{55}} & \textcolor{myred}{\ding{55}} & \textcolor{myred}{\ding{55}} & \textcolor{myred}{\ding{55}} & \textcolor{myred}{\ding{55}} & \textcolor{myred}{\ding{55}}\\
Lifelong Learning & \textcolor{myred}{\ding{55}} & \textcolor{myred}{\ding{55}} & \textcolor{mygreen}{\checkmark} & \textcolor{myred}{\ding{55}} & \textcolor{myred}{\ding{55}} & \textcolor{myred}{\ding{55}} & \textcolor{myred}{\ding{55}} \\
Model Editing & \textcolor{myred}{\ding{55}} & \textcolor{myred}{\ding{55}} & \textcolor{mygreen}{\checkmark} & \textcolor{mygreen}{\checkmark} & \textcolor{myred}{\ding{55}} & \textcolor{myred}{\ding{55}} & \textcolor{myred}{\ding{55}} \\
\midrule
\textbf{Self-evolving Agents} & 
\textcolor{mygreen}{\checkmark} &
\textcolor{mygreen}{\checkmark} &
\textcolor{mygreen}{\checkmark} & \textcolor{mygreen}{\checkmark} & \textcolor{mygreen}{\checkmark} & \textcolor{mygreen}{\checkmark} & \textcolor{mygreen}{\checkmark}\\
\bottomrule
\end{tabular}
}
\label{tab:compare-paradigm}
\end{table}



\paragraph{Positioning Self-Evolving Agents}

To clarify the relationships among these paradigms and to motivate the role of self-evolving agents, we examine them through two complementary perspectives: a \emph{problem-setting} lens and a \emph{solution-paradigm} lens. This distinction clarifies the basis of each paradigm - whether it emerges from constraints and challenges inherent to the learning setting, or from methodological proposals for how the model or agent itself can be updated.

\begin{itemize}


\item \textbf{Problem-setting view.}
Curriculum learning and lifelong learning arise from concrete learning problems. Curriculum learning addresses how to structure training examples of varying difficulty so a model can handle complex samples more effectively; lifelong learning focuses on acquiring new abilities over time while mitigating catastrophic forgetting. These paradigms are therefore driven by the \emph{problems} they aim to solve and primarily specify how experience is organized for the learner, rather than how the agent itself may adapt beyond parameter updates.

\item \textbf{Solution-paradigm view.}
Model editing and self-evolving agents, in contrast, originate as \emph{solutions}: they propose mechanisms for updating or modifying a system. Model editing provides targeted procedures—typically localized parameter adjustments—to correct or insert knowledge. Self-evolving agents generalize this idea by treating adaptation as a first-class capability, allowing not only parameter updates but also changes to runtime context, memory, tools, and workflow structures, driven by the agent’s own trajectories and feedback signals.

\end{itemize}

Viewed through this two-lens framework, curriculum and lifelong learning are anchored in the nature of the learning \emph{problems} they address, whereas model editing and self-evolving agents are defined by the \emph{methods} they provide for effecting change. Self-evolving agents thus represent a system-level solution paradigm: they include parameter-level editing as one update pathway while enabling broader, persistent, and interaction-driven evolution across multiple components of an agent.




